\newpage
\pagestyle{fancy}
\lhead{\textsc{Chapitre 3. Choix Technologiques et Conception}}
\renewcommand{\headrulewidth}{0.5pt}
\renewcommand{\footrulewidth}{0.5pt}
\chapter{Choix Technologiques et Conception}
\setlength{\headheight}{27.06pt}
\label{ch3}
\minitoc
\newpage
\section*{Introduction}
\addcontentsline{toc}{section}{Introduction}
Ce chapitre présente les choix technologiques retenus pour la réalisation de la plateforme HumatiQ, ainsi que la conception architecturale du système. Nous commençons par justifier le choix de chaque technologie selon les exigences du projet, en couvrant le frontend, le backend, les bases de données et les composantes d'intelligence artificielle. Nous exposons ensuite l'architecture globale du système et les principaux diagrammes de conception.
% ================================================================
\section{Choix technologiques}
% ================================================================
\subsection{Frontend : React.js}
React.js\footnote{\url{https://react.dev}} est une bibliothèque JavaScript développée et maintenue par Meta, dédiée à la construction d'interfaces utilisateur dynamiques sous forme de composants réutilisables. Associé à Vite\footnote{\url{https://vitejs.dev}} comme outil de build, il offre un cycle de développement rapide grâce au \textit{Hot Module Replacement} (HMR) et à un bundling optimisé.
\noindent Les raisons qui ont motivé ce choix sont les suivantes :
\begin{itemize}
    \item \textbf{Architecture composants} : facilite la modularisation des interfaces complexes (pipeline Kanban, tableau de bord analytique, salle d'entretien vidéo).
    \item \textbf{Écosystème riche} : intégration native avec des bibliothèques spécialisées telles que \texttt{recharts} pour les graphiques, \texttt{framer-motion} pour les animations fluides, \texttt{react-webcam} pour la capture vidéo et \texttt{@mediapipe/selfie\_segmentation} pour la segmentation d'arrière-plan en temps réel lors des entretiens vidéo.
    \item \textbf{Styling moderne} : Tailwind CSS v4\footnote{\url{https://tailwindcss.com}} est utilisé comme framework de styles utilitaires, permettant un design cohérent, responsive et maintenable sans quitter le fichier JSX.
    \item \textbf{Communication temps réel} : le frontend communique avec le backend via l'API WebSocket native du navigateur pour le signaling WebRTC et les notifications en temps réel, sans dépendance à une bibliothèque tierce.
    \item \textbf{Performance} : le DOM virtuel de React minimise les re-rendus inutiles, garantissant une fluidité optimale pour des vues intensives en données.
    \item \textbf{SPA (\textit{Single Page Application})} : navigation sans rechargement de page, essentielle pour l'expérience utilisateur de la plateforme.
\end{itemize}
\noindent Le thème visuel adopté, baptisé \textit{Architectural Curator}, repose sur un design sombre avec effets de glassmorphisme et accents dorés, offrant une identité visuelle distincte et professionnelle.
% ----------------------------------------------------------------
\subsection{Backend : FastAPI}
FastAPI\footnote{\url{https://fastapi.tiangolo.com}} est un framework Python moderne orienté hautes performances, basé sur les standards OpenAPI et JSON Schema. Il tire parti des annotations de type Python et de la programmation asynchrone (\texttt{async/await}) pour exposer des API REST robustes et documentées automatiquement.
\noindent Les atouts déterminants de FastAPI pour HumatiQ sont :
\begin{itemize}
    \item \textbf{Performances asynchrones} : gestion efficace des I/O concurrentes, indispensable pour les endpoints intensifs en IA (scoring, génération de quiz, analyse vidéo). Le driver asynchrone \texttt{motor} est utilisé pour toutes les interactions MongoDB.
    \item \textbf{Documentation automatique} : génération native de la documentation interactive Swagger UI et ReDoc à partir des schémas Pydantic, facilitant l'intégration frontend/backend.
    \item \textbf{Validation des données} : Pydantic assure une validation stricte et typée de toutes les entrées/sorties de l'API.
    \item \textbf{WebSocket natif} : support intégré des connexions WebSocket pour le signaling WebRTC et les notifications en temps réel.
    \item \textbf{Écosystème Python IA} : interopérabilité native avec les bibliothèques d'intelligence artificielle (\texttt{transformers}, \texttt{torch}, \texttt{torchvision}, \texttt{opencv-python}).
\end{itemize}
% ----------------------------------------------------------------
\subsection{Bases de données}
La nature hétérogène des données manipulées par HumatiQ — profils candidats, documents CV, offres d'emploi, résultats d'entretiens — a nécessité une réflexion approfondie sur le choix du système de gestion de bases de données.
\subsubsection{Approche Relationnelle}
Les bases de données relationnelles (SQL) reposent sur un modèle tabulaire structuré, avec des relations explicites entre entités, l'intégrité référentielle et le support des transactions ACID. Elles sont particulièrement adaptées aux données fortement structurées et aux requêtes analytiques complexes.
\paragraph{Supabase}
Supabase\footnote{\url{https://supabase.com}} est une plateforme open-source de \textit{Backend-as-a-Service} construite sur PostgreSQL. Elle offre une authentification intégrée (OAuth2, OTP, JWT), un stockage de fichiers, des fonctions Edge et une API REST/GraphQL générée automatiquement.
Dans le contexte de HumatiQ, Supabase est utilisé pour la gestion des identités et des sessions utilisateurs, en tirant parti de son module d'authentification éprouvé et de sa compatibilité OAuth2 avec Google. Le client officiel \texttt{@supabase/supabase-js} est également intégré côté frontend pour la gestion des tokens JWT.
\subsubsection{Approche NoSQL}
Les bases de données NoSQL adoptent des modèles de stockage flexibles (document, clé-valeur, graphe, colonne) et sont conçues pour absorber des volumes de données hétérogènes et évolutifs, sans contrainte de schéma rigide.
\paragraph{Comparaison SQL vs NoSQL}
\begin{table}[H]
\centering
\renewcommand{\arraystretch}{1.4}
\resizebox{\textwidth}{!}{%
\begin{tabular}{|l|c|c|}
\hline
\rowcolor{gray!80}
\textcolor{white}{\textbf{Critère}} &
\textcolor{white}{\textbf{SQL (PostgreSQL / Supabase)}} &
\textcolor{white}{\textbf{NoSQL (MongoDB)}} \\
\hline
Schéma & Rigide, défini à l'avance & Flexible, évolutif \\
\hline
Requêtes complexes & Très puissantes (JOIN, agrégats) & Limitées sans pipeline d'agrégation \\
\hline
Scalabilité horizontale & Difficile & Native \\
\hline
Données hétérogènes & Mal adaptée & Idéale \\
\hline
Transactions ACID & Natives & Supportées (depuis v4.0) \\
\hline
Stockage de fichiers & Extension (GridFS, objets) & GridFS natif \\
\hline
Cas d'usage & Authentification, données structurées & Profils candidats, CVs, entretiens \\
\hline
\end{tabular}%
}
\caption{Comparaison SQL vs NoSQL pour HumatiQ}
\label{tab:sql_nosql}
\end{table}
\paragraph{MongoDB}
MongoDB\footnote{\url{https://www.mongodb.com}} est le système de base de données orienté documents le plus répandu dans l'écosystème NoSQL. Chaque enregistrement est stocké sous forme de document BSON (\textit{Binary JSON}), permettant une représentation native des structures hiérarchiques et des listes imbriquées.
\noindent Les raisons de son adoption pour HumatiQ sont :
\begin{itemize}
    \item \textbf{Flexibilité des schémas} : les profils candidats varient fortement d'un utilisateur à l'autre (nombre d'expériences, certifications, compétences). MongoDB absorbe cette hétérogénéité sans migration de schéma.
    \item \textbf{GridFS} : module natif de stockage de fichiers binaires volumineux (CVs PDF, documents de quiz), évitant de recourir à un service de stockage externe.
    \item \textbf{Pipeline d'agrégation} : puissant moteur de traitement analytique pour les tableaux de bord KPI.
    \item \textbf{Atlas Vector Search} : support natif de la recherche vectorielle via l'opérateur \texttt{\$vectorSearch}, essentiel pour le matching IA basé sur les embeddings de profils candidats.
\end{itemize}
\subsubsection{Choix : Architecture Hybride}
Face aux forces complémentaires des deux paradigmes, nous avons opté pour une \textbf{architecture hybride} :
\begin{itemize}
    \item \textbf{Supabase (PostgreSQL)} gère l'authentification, la gestion des sessions et des tokens JWT, les données d'entreprise structurées et les relations entre rôles.
    \item \textbf{MongoDB} stocke les profils candidats, les offres d'emploi, les candidatures, les quiz, les transcriptions d'entretiens et les fichiers binaires via GridFS.
\end{itemize}
Cette séparation des responsabilités permet d'exploiter les forces de chaque système tout en maintenant une architecture cohérente et maintenable.
% ----------------------------------------------------------------
\subsection{Intelligence Artificielle}
L'intelligence artificielle constitue le cœur différenciateur de HumatiQ. Plusieurs approches ont été évaluées avant d'arrêter le choix définitif.
\subsubsection{Modèles Locaux : Ollama}
Ollama\footnote{\url{https://ollama.com}} est un outil permettant d'exécuter des modèles de langage de grande taille (LLM) en local, sans dépendance à une API cloud. Il supporte des modèles open-source tels que LLaMA 3, Mistral et Qwen.
\noindent \textbf{Avantages} : confidentialité totale des données, absence de coûts d'inférence, latence maîtrisée.\\
\noindent \textbf{Inconvénients} : nécessite des ressources GPU importantes, performances pouvant être inférieures aux modèles cloud pour des tâches de raisonnement très complexes.
\subsubsection{Modèles Cloud : HuggingFace}
HuggingFace\footnote{\url{https://huggingface.co}} est la plateforme de référence pour l'hébergement et le déploiement de modèles de machine learning. Elle offre via son \textit{Inference API} un accès à des milliers de modèles pré-entraînés (Qwen2.5, Whisper, etc.) à travers des appels REST standardisés. Pour HumatiQ, cette API est également utilisée pour l'analyse post-entretien avec le modèle \texttt{Qwen/Qwen2.5-72B-Instruct}.
\noindent \textbf{Avantages} : accès à des modèles spécialisés de pointe, sans infrastructure GPU.\\
\noindent \textbf{Inconvénients} : dépendance réseau, coût variable selon le volume d'inférence.
\subsubsection{Modèles entraînés en interne}
Pour des besoins très spécifiques à HumatiQ — notamment la détection d'émotions faciales en temps réel pendant les entretiens — des modèles ont été entraînés en interne sur des données annotées. Cette approche garantit une précision supérieure sur le domaine cible, au prix d'un effort de collecte et d'entraînement plus important.
\subsubsection{Comparaison des approches IA}
\begin{table}[H]
\centering
\renewcommand{\arraystretch}{1.4}
\resizebox{\textwidth}{!}{%
\begin{tabular}{|l|c|c|c|}
\hline
\rowcolor{gray!80}
\textcolor{white}{\textbf{Critère}} &
\textcolor{white}{\textbf{Ollama (local)}} &
\textcolor{white}{\textbf{HuggingFace (cloud)}} &
\textcolor{white}{\textbf{Modèles entraînés}} \\
\hline
Confidentialité des données & \cmark\cmark & \xmark & \cmark\cmark \\
\hline
Coût d'inférence & Nul & Variable & Nul (après entraînement) \\
\hline
Qualité sur tâches générales & Élevée & Très élevée & Limitée au domaine \\
\hline
Qualité sur domaine métier & Moyenne & Moyenne & Très élevée \\
\hline
Dépendance infrastructure & GPU local & Réseau & GPU local \\
\hline
Facilité de déploiement & Moyenne & Élevée & Faible \\
\hline
\end{tabular}%
}
\caption{Comparaison des approches IA pour HumatiQ}
\label{tab:comparaison_ia}
\end{table}
\subsubsection{Choix : Mix Technologique d'IA}
Afin d'optimiser qualité, coût et confidentialité, HumatiQ adopte un \textbf{mix technologique d'IA} configurable par variables d'environnement :
\begin{itemize}
    \item \textbf{LLM via Ollama (local) ou API cloud (OpenAI / HuggingFace)} : scoring multicritères des CVs, génération de quiz par RAG. Ollama est le fournisseur par défaut ; OpenAI et HuggingFace constituent des alternatives cloud activables via configuration. Cette architecture flexible permet de basculer entre modes local et cloud sans modification du code.
    \item \textbf{Embeddings \texttt{nomic-embed-text}} : modèle open-source d'embeddings vectoriels (768 dimensions) exécuté via Ollama, utilisé pour le pipeline de recherche sémantique des candidats et pour la génération des quiz par RAG.
    \item \textbf{Modèle FER entraîné en interne} : réseau de neurones à architecture ResNet entraîné sur le dataset FER-2013 pour la détection des 7 émotions faciales en temps réel pendant les entretiens vidéo. La détection des visages s'appuie sur le classifieur Haar Cascade d'OpenCV.
    \item \textbf{Analyse post-entretien via HuggingFace} : le modèle \texttt{Qwen/Qwen2.5-72B-Instruct} est appelé via l'Inference API de HuggingFace pour agréger la transcription, les données émotionnelles et les notes du recruteur en un rapport structuré (résumé, points forts, axes d'amélioration, score global).
    \item \textbf{Whisper (via HuggingFace Transformers)} : modèle de transcription parole-texte chargé via la bibliothèque \texttt{transformers} d'HuggingFace, utilisé pour la génération automatique des transcriptions d'entretiens.
\end{itemize}
% ================================================================
\section{Architecture et Conception du système}
% ================================================================
\subsection{Architecture Globale}
L'architecture de HumatiQ repose sur un modèle \textbf{client-serveur multicouches}, séparant clairement les responsabilités entre le frontend, le backend API, la couche IA et les systèmes de persistance.
\begin{figure}[H]
    \centering
    \includegraphics[width=1\linewidth]{images/architecture_globale.png}
    \caption{Architecture globale de HumatiQ}
    \label{fig:architecture_globale}
\end{figure}
\noindent Les quatre couches principales de l'architecture sont :
\begin{itemize}
    \item \textbf{Couche Présentation (Frontend)} : application React.js/Vite communicant avec le backend via des requêtes HTTP REST et des connexions WebSocket natives pour le temps réel (notifications, signaling WebRTC des entretiens vidéo).
    \item \textbf{Couche API (Backend)} : serveur FastAPI exposant l'ensemble des endpoints REST, gérant l'authentification JWT, le routage des requêtes et l'orchestration des services IA. Le driver asynchrone \texttt{motor} assure les interactions non-bloquantes avec MongoDB.
    \item \textbf{Couche IA} : ensemble de modules Python spécialisés (scoring LLM, embeddings via Ollama, FER/ResNet, Whisper, analyse post-entretien via HuggingFace) exposés comme services internes consommés par le backend.
    \item \textbf{Couche Persistance} : architecture hybride Supabase (authentification, données structurées) et MongoDB (profils, offres, candidatures, fichiers GridFS, vecteurs d'embeddings).
\end{itemize}
\noindent Le tableau suivant récapitule l'ensemble du stack technologique retenu :
\begin{table}[H]
\centering
\renewcommand{\arraystretch}{1.4}
\resizebox{\textwidth}{!}{%
\begin{tabular}{|l|l|p{8cm}|}
\hline
\rowcolor{gray!80}
\textcolor{white}{\textbf{Couche}} &
\textcolor{white}{\textbf{Technologie}} &
\textcolor{white}{\textbf{Rôle et justification}} \\
\hline
Backend API & FastAPI (Python 3.11+) & Framework asynchrone haute performance. Génération automatique de documentation OpenAPI. Idéal pour des endpoints IA intensifs. \\
\hline
Base de données & MongoDB (Atlas compatible) & Stockage documentaire flexible adapté aux profils candidats hétérogènes. GridFS pour les fichiers CV. Atlas Vector Search pour le matching sémantique. \\
\hline
Authentification & Supabase + JWT + OAuth2 & Gestion sécurisée des identités. OAuth2 Google pour l'onboarding simplifié des candidats et RH. \\
\hline
Temps réel & WebSocket (natif Python + navigateur) & Signaling WebRTC pour les entretiens vidéo, notifications in-app en temps réel. \\
\hline
Frontend & React.js + Vite + Tailwind CSS v4 & SPA performante. Vite assure un bundling rapide. Tailwind CSS pour le design utilitaire. Thème sombre glassmorphisme avec accents dorés. \\
\hline
Animations & Framer Motion & Animations d'interface fluides et transitions entre vues. \\
\hline
Vidéo P2P & WebRTC + MediaPipe & Entretiens vidéo peer-to-peer. MediaPipe Selfie Segmentation pour la suppression d'arrière-plan en temps réel. \\
\hline
IA — Scoring & LLM (Ollama / OpenAI / HuggingFace) + Embeddings & Scoring multicritères des CVs, génération de quiz contextuels par RAG. Fournisseur configurable via variables d'environnement. \\
\hline
IA — Émotion & ResNet (Deep Learning) & Analyse faciale en temps réel (7 émotions). Architecture ResNet entraînée sur FER-2013, détection de visages via Haar Cascade (OpenCV). \\
\hline
IA — Analyse entretien & Qwen2.5-72B-Instruct (HuggingFace) & Génération du rapport post-entretien à partir de la transcription et des données émotionnelles. \\
\hline
IA — Transcription & Whisper (HuggingFace Transformers) & Transcription parole-texte pour les entretiens vidéo, chargé via la bibliothèque \texttt{transformers}. \\
\hline
Embeddings & \texttt{nomic-embed-text} (Ollama) & Modèle open-source d'embeddings vectoriels 768 dimensions pour le pipeline RAG des quiz et le matching candidats. \\
\hline
OCR / Extraction & PyMuPDF + Tesseract & Extraction de texte depuis PDF, images et documents numérisés. \\
\hline
Agenda & Google Calendar API & Synchronisation automatique des entretiens planifiés via OAuth2. \\
\hline
Email & SMTP (configurable) & Envoi de toutes les communications automatiques (invitations, rappels, résultats). \\
\hline
\end{tabular}%
}
\caption{Stack technologique de HumatiQ}
\label{tab:stack}
\end{table}
% ----------------------------------------------------------------
\subsection{Diagramme de classes}
Le diagramme de classes présenté en figure~\ref{fig:classes} modélise les entités principales de la plateforme HumatiQ et leurs relations. Les treize collections MongoDB identifiées sont représentées avec leurs attributs essentiels et les associations qui les relient.
\begin{figure}[H]
    \centering
    \includegraphics[width=1\linewidth]{images/diagramme_classes.png}
    \caption{Diagramme de classes de HumatiQ}
    \label{fig:classes}
\end{figure}
\noindent Les entités principales du modèle sont :
\begin{itemize}
    \item \textbf{User} : entité centrale regroupant les attributs communs à tous les utilisateurs (identifiant Supabase UUID, email, rôle, statut d'authentification 2FA). Spécialisée en \textit{Candidate}, \textit{HRUser} et \textit{SuperAdmin}.
    \item \textbf{Company} : entreprise cliente de la plateforme, associée à ses départements et à ses offres d'emploi.
    \item \textbf{JobOffer} : offre d'emploi publiée par un recruteur, contenant les paramètres de présélection, la date limite, la configuration d'automatisation IA et les paramètres du pipeline de quiz.
    \item \textbf{Application} : candidature d'un profil à une offre, portant le statut Kanban courant, le score IA, le score quiz et les évaluations associées.
    \item \textbf{CandidateProfile} : profil enrichi du candidat, incluant ses expériences, formations, certifications et le vecteur d'embedding (768 dimensions) de son profil.
    \item \textbf{Quiz} : évaluation générée par IA à partir de documents via RAG, composée de questions QCM ou Vrai/Faux, publiée à un candidat avec une date limite et un minuteur.
    \item \textbf{Interview} : entretien vidéo planifié, associant un recruteur, un candidat, un créneau horaire, un historique d'émotions, une transcription horodatée et un rapport post-entretien généré par IA.
\end{itemize}
% ----------------------------------------------------------------
\subsection{Diagrammes de séquences}
Les diagrammes de séquences illustrent les interactions dynamiques entre les acteurs et les composantes du système pour les scénarios fonctionnels clés. Nous présentons ci-après les séquences relatives au scoring IA des candidatures et à la conduite d'un entretien vidéo.
\subsubsection{Séquence : Scoring IA d'une candidature}
\begin{figure}[H]
    \centering
    \includegraphics[width=1\linewidth]{images/seq_scoring.png}
    \caption{Diagramme de séquence — Scoring IA d'une candidature}
    \label{fig:seq_scoring}
\end{figure}
\noindent Le processus de scoring se déroule selon les étapes suivantes :
\begin{enumerate}
    \item Le recruteur déclenche le scoring depuis le pipeline de candidatures.
    \item Le backend construit un texte sémantique structuré à partir du profil candidat stocké en MongoDB (titre, compétences, expériences, formations, certifications) et génère son vecteur d'embedding via le modèle \texttt{nomic-embed-text} exécuté sur Ollama.
    \item La description du poste est également encodée en vecteur par le même modèle.
    \item MongoDB Atlas Vector Search (\texttt{\$vectorSearch}) réalise la recherche par similarité entre le vecteur de l'offre et les vecteurs des profils candidats, retournant un score de pertinence normalisé.
    \item Le LLM (Ollama par défaut, OpenAI ou HuggingFace en mode cloud) réalise une analyse multicritères (compétences : 40 pts, expérience : 40 pts, formation : 20 pts) et retourne un score numérique avec justification textuelle au format JSON.
    \item Le score final est enregistré dans la collection \textit{applications} de MongoDB et affiché dans l'interface Kanban.
\end{enumerate}
\subsubsection{Séquence : Conduite d'un entretien vidéo}
\begin{figure}[H]
    \centering
    \includegraphics[width=1\linewidth]{images/seq_entretien.png}
    \caption{Diagramme de séquence — Entretien vidéo en temps réel}
    \label{fig:seq_entretien}
\end{figure}
\noindent Le déroulement d'un entretien vidéo implique les flux simultanés suivants :
\begin{enumerate}
    \item Établissement de la connexion WebRTC via le serveur de signaling WebSocket intégré à FastAPI.
    \item Côté frontend, MediaPipe Selfie Segmentation applique la suppression d'arrière-plan en temps réel sur le flux vidéo du candidat avant transmission.
    \item Les frames vidéo sont envoyées au backend via WebSocket ; le modèle ResNet (FER) analyse chaque frame pour détecter et horodater les émotions faciales (7 classes : angry, disgust, fear, happy, neutral, sad, surprise).
    \item La transcription de l'audio est assurée par Whisper (via HuggingFace Transformers) et les entrées sont enregistrées avec horodatage dans la collection \textit{interviews}.
    \item À la clôture de l'entretien, le modèle \texttt{Qwen/Qwen2.5-72B-Instruct} (HuggingFace Inference API) agrège la transcription complète et l'historique émotionnel pour générer un rapport structuré (résumé, points forts, axes d'amélioration, score global).
\end{enumerate}
% ================================================================
\section*{Conclusion}
\addcontentsline{toc}{section}{Conclusion}
% ================================================================
Ce chapitre a permis de poser les fondations techniques de la plateforme HumatiQ. Nous avons justifié les choix technologiques retenus pour chaque couche du système — React.js/Vite/Tailwind CSS pour le frontend, FastAPI avec \texttt{motor} pour le backend, une architecture hybride Supabase/MongoDB pour la persistance, et un mix technologique d'IA combinant Ollama pour les modèles locaux, HuggingFace pour l'analyse post-entretien, et des modèles entraînés en interne pour la détection d'émotions. Nous avons ensuite présenté l'architecture globale du système ainsi que les diagrammes de classes et de séquences modélisant les interactions clés. Le chapitre suivant est consacré au développement du premier sprint, couvrant l'authentification et la gestion des utilisateurs.
